\section{Learning the groups}\label{sec:algorithm}
The algorithm turns the data into modes (\cref{sec:modes}), clusters the modes onto the curves of factors within each
sample of the field (grouping I, \cref{sec:group1}) inside a class of operators chosen by the data
(\cref{sec:class}), clusters the factors across samples (grouping II, \cref{sec:group2}), and labels the result
(\cref{sec:labels}). The complete algorithm, with every parameter and its default, is in \SMalgo.

\subsection{From data to modes}\label{sec:modes}
\paragraph{Field samples} The input is a set of noisy observations $y_i=u(t_i,x_i)+\varepsilon_i$,
$\varepsilon_i\sim\Normal(0,\noise^2)$ (here and below $\Normal(m,C)$ is a Gaussian with mean $m$ and covariance $C$), on a window $[0,T]\times[0,L)$ periodic in $x$, or a noisy field on a regular grid.
Scattered observations are interpolated by a Gaussian process, a probability distribution over functions. The data are
modelled as $y=\bar y+f+\varepsilon$ with a constant $\bar y$ and
\begin{equation}\label{eq:gp}
  f\sim\mathrm{GP}\Big(0,\ \sum_{r=1}^R\eta_r^2\,k_t^{(r)}(t,t')\,k_x^{(r)}(x,x')\Big),
\end{equation}
where each $k^{(r)}$ is a Mat\'ern-$5/2$ kernel with its own length scale ($k_x$ periodic in $x$) and $\eta_r$ is its
amplitude; several products are used because a transported profile and a diffusing one do not share a length scale.
Conditioned on the data, the process is again a Gaussian distribution over functions: a mean function, and a
covariance that says how far plausible fields can deviate from it. It contains infinitely many functions, which cannot be
enumerated, so the algorithm represents it by $\Nsamp$ random draws $u^{(q)}$, each a plausible field consistent with the
data, and runs every later step on each of them; the spread of the answers over the draws is the uncertainty. The
hyperparameters (amplitudes, length scales, noise level) are themselves unknown: they are drawn from an approximation of
their posterior and reweighted by how well they explain the data (importance sampling). The posterior is therefore a
weighted mixture of Gaussian processes, not a single one, and each draw $u^{(q)}$ carries the weight $\pi_q$ of the
hyperparameters it was drawn with. The posterior mean is deliberately not used. The reason is structural: the maps applied next are not continuous
(factors of a numerical symbol jump where it becomes reducible), so the mean of the outputs differs from the output of
the mean, and the mean, a Wiener filter, shrinks exactly the fast-decaying high-wavenumber modes that carry the
diffusive factors. On gridded input the noisy grid itself is used, and parametric-bootstrap refits provide the spread.

\paragraph{The usable band} Not every wavenumber carries information about the symbol. Three limits apply: spatial
aliasing ($\abs{\kappa}\le\pi/\Delta x$); temporal aliasing along a branch, since a rate is read from samples spaced by
$\Delta t$ only modulo $2\pi i/\Delta t$; and, usually the binding one, noise together with decay. After the spatial Fourier
transform one wavenumber bin carries noise of standard deviation $\sigbin=\noise/\sqrt{n_x}$, and a mode is used only
while its amplitude exceeds $z\,\sigbin$ ($z=3$). A diffusive mode of rate $-\nu\kappa^2$ falls below this threshold after
a time of order
\[
  t_\kappa\approx\frac{\log\big(\abs{A_\kappa}/z\sigbin\big)}{\nu\kappa^2},
\]
so the high wavenumbers of a dissipative group are lost
first. This is a property of the data, not of the method. The band $\mathcal K$ is the set of wavenumbers that pass all
three limits.

\paragraph{How many wavenumbers does a group need?} A factor is a curve, and a curve of degree $d$ is not fixed by too
few of its points. The band must therefore hold enough distinct wavenumbers on each group:

\begin{proposition}[Counting condition]\label{prop:counting}
Let $Q=\tau-\Lambda(\xi)$ with $\Lambda$ a polynomial of degree at most $d$. Its coefficients are determined by the rates
$\lambda(\kappa_\ell)=\Lambda(i\kappa_\ell)$ at $d+1$ distinct wavenumbers, and not by fewer; the coefficient error is at
most $\norm{V^{-1}}\max_\ell\abs{\delta\lambda_\ell}$ with $V=((i\kappa_\ell)^k)$ the Vandermonde matrix, whose inverse grows
like $h_\kappa^{-d}$ when the nodes cluster within a spread $h_\kappa$~\cite{Gau62}. For real coefficients
$\lfloor d/2\rfloor+1$ distinct values of $\abs{\kappa}$ suffice.
\end{proposition}

So the whole band is used, not the cleanest few bins: spreading the wavenumbers is what keeps $V$ well conditioned.

\paragraph{Rates by the matrix pencil} For each sample and each $\kappa\in\mathcal K$, the spatial Fourier transform gives
a time series that, for a constant-coefficient equation, is a sum of exponentials, one per root of the symbol at that
wavenumber:
\begin{equation}\label{eq:expsum}
  \hat u(t,\kappa)=\sum_{p=1}^{r(\kappa)}A_p\,e^{\lambda_p(\kappa)t},\qquad Q_{j(p)}\big(\lambda_p,i\kappa\big)=0 .
\end{equation}
The rates are the eigenvalues of $\partial_t$ on the space spanned by the series and its shifts, and the matrix pencil
method computes them~\cite{HuaSar90}. From the Hankel matrix $H_{\ell k}=\hat u(t_{\ell+k},\kappa)$ take $Y_0$ ($H$ without
its last column) and $Y_1$ ($H$ without its first), both projected on the $r$ leading singular directions. Advancing a
series by one step multiplies the term $A_pe^{\lambda_pt}$ by $\mu_p=e^{\lambda_p\Delta t}$, so
\begin{equation}\label{eq:pencil}
  \lambda_p=\tfrac{1}{\Delta t}\log\mu_p,\qquad \mu_p\in\operatorname{eig}\big(Y_0^{+}Y_1\big).
\end{equation}
The order $r$ is chosen per bin by BIC, and each rate receives a standard deviation $\sigma_{\lambda,p}$ from the
Cram\'er--Rao bound of model \eqref{eq:expsum}; weak, fast-decaying or nearly coincident modes receive large ones.
This step does \emph{not} fit the equation: it returns a cloud of points $(\kappa,\lambda_p(\kappa))$ with error bars, the
noisy version of \cref{fig:branches}. Two rates closer than about $4/T$ cannot be resolved in a window of length $T$; such
bins are flagged as confluent.

\subsection{Grouping I: modes onto curves}\label{sec:group1}
The obvious route, fitting one polynomial through all points and factoring it, is ill-posed: a polynomial with noisy
coefficients is almost never reducible, and when it is, its factors move discontinuously with the coefficients (in our
tests this route split the irreducible Klein--Gordon and telegraph symbols at noise $10^{-4}$). We therefore group and fit
at the same time. Given a structure $\Mfit$ (the number $N$ of factors and their degrees, \cref{sec:class}), we solve the
\emph{$N$-varieties problem}, a version of $k$-means in which the cluster centres are curves instead of points: assign
every mode to the factor whose variety is nearest, refit each factor to the modes assigned to it, and repeat until the
assignment stops changing. The distance of mode $p$ to factor $Q_j$ is its first-order distance to the variety in units
of its own uncertainty,
\begin{equation}\label{eq:dist}
  r_j(p)=\frac{\abs{Q_j(\lambda_p,i\kappa_p)}}{\sigma_{\lambda,p}\,\abs{\partial_\tau Q_j(\lambda_p,i\kappa_p)}},
\end{equation}
and each mode costs $\min(r_j(p)^2,9)$, so an isolated outlier pays a fixed penalty instead of dragging a curve away.
This is the grouping of \cref{thm:factor} learned directly from the spectrum.

\begin{figure}[tbp]
\centering
\includegraphics[width=\linewidth]{fig_directions}
\caption{Grouping I on the running example ($64\times128$ grid, $5\%$ noise). Points: modes found by the matrix pencil,
coloured by the factor they were assigned to. Dashed: the true branches $\lambda=-ic\kappa$ and $\lambda=-\nu\kappa^2$.
Thin lines: the factors fitted to $32$ posterior samples; they bundle where there are modes and fan out where there are
none. (a) Oscillation $\operatorname{Im}\lambda$: transport is the sloped line, diffusion the flat one. (b) Damping
$\operatorname{Re}\lambda$: transport is flat at zero, diffusion the downward parabola.}\label{fig:directions}
\end{figure}

\Cref{fig:directions} shows the result on the running example. Each factor is a direction in the plane of rates, and the
modes of one group line up along it. Once modes are assigned, the component of group $j$ is the sum of its modes,
\begin{equation}\label{eq:components}
  u_j(t,x)=\sum_{p:\,j(p)=j}A_p\,e^{\lambda_p(t-t_0)+i\kappa_px},
\end{equation}
the discrete counterpart of the projector \eqref{eq:lagrange}. Indeed, on the sum of exponentials \eqref{eq:expsum} each
factor $\partial_t-\Lambda_k(i\kappa)$ of that multiplier annihilates the modes of branch $k$ and multiplies a mode of
branch $j$ by $\Lambda_j-\Lambda_k$, which the denominator cancels: the projector keeps exactly the modes of group $j$,
and this is \eqref{eq:components}. A mode within distance $3$ of two varieties lies in the overlap $\deltaP$ (for the running example, $\kappa=0$) and is split
equally between the groups: among all splits allowed by \cref{thm:factor} this one puts the least energy into the
ambiguity, and it coincides with the posterior mean of \cref{prop:bayes-angle}. The energy shares
$w_j=\norm{u_j}/\sum_k\norm{u_k}$ follow. A group can itself split further (two conjugate branches of the beam equation are
first found as one factor of degree two in $\tau$); each component is therefore analysed again, and a split is kept only
if it leaves the kernel unchanged or improves the evidence of \cref{sec:labels}.

\subsection{Choosing the class of operators}\label{sec:class}
A structure $\Mfit=((m_1,d_1),\dots,(m_N,d_N))$ fixes $N$ and, for each factor, its degree $m_j$ in $\tau$ and $d_j$ in
$\xi$; the running example is $((1,1),(1,2))$. The class $C(m,d)$ holds all structures with $\sum_jm_j\le m$ and
$d_j\le d$. Grouping I is run for every structure, and structures are compared by
\begin{align}
  \mathrm{BIC}(\Mfit)&=2n_m\log\max\!\Big(\frac{\chi^2_{\mathrm{modes}}}{2n_m},\,1\Big)+n_{\mathrm{par}}(\Mfit)\log(2n_m),\label{eq:bic}\\
  P(\Mfit\mid\text{data})&\propto e^{-\mathrm{BIC}(\Mfit)/2},\label{eq:pstruct}
\end{align}
where $n_m$ is the number of fitted modes, $n_{\mathrm{par}}$ the number of coefficients of the structure, and
\[
  \chi^2_{\mathrm{modes}}=\sum_{p=1}^{n_m}\min\big(r_{j(p)}(p)^2,\,9\big)
\]
is the misfit of the modes: the squared distance~\eqref{eq:dist} of each mode to the curve of its factor, in units of its
own standard deviation, summed over the modes. On the data it is one number. If the structure is right, the same number
computed over repeated noise realizations would follow approximately one half of a chi-squared distribution with
$2n_m-n_{\mathrm{par}}$ degrees of freedom (each complex rate gives two real data), so $\chi^2_{\mathrm{modes}}/(2n_m)\approx\frac12$
is what noise alone produces. The first term of \eqref{eq:bic} penalizes misfit beyond that level, and the floor
$\max(\cdot,1)$ means that a fit better than the noise is never rewarded; the second term penalizes coefficients, and
\eqref{eq:pstruct} turns the trade-off into a probability over structures. This is Occam's razor made quantitative:
among structures that explain the data equally well, the one with fewer coefficients wins. It is also the preference
for simple laws voiced by Einstein, who asked that the basic elements of a theory be as simple and as few as possible
without giving up the adequate representation of experience~\cite{Ein34}, and by Dirac, who held that the fundamental
equations of physics should have mathematical beauty~\cite{Dir63}; the criterion keeps simplicity subordinate to the
data and never trades fit for it. The class is grown from $C(3,4)$ (90 structures) up to $C(4,8)$ (1353) only when the data demand
it, i.e.\ when most of the posterior mass is on structures that fail the tests of \cref{sec:labels}, and a larger class is
kept only if it also predicts the second half of the time window from the first. This is the only regularization of the
method, and it acts on the class, by Occam's razor, not on the coefficients: by \cref{sec:factor} a coefficient penalty
cannot remove the interaction of terms, and by \cref{thm:hm} it can push the error of the split below its floor only by
shrinking the components, which the next layer does explicitly and reports.

\subsection{Grouping II: factors across samples}\label{sec:group2}
Grouping I gives one set of factors per field sample. Two questions remain, one about solutions and one about factors:
given the factors, how well are the components $u_j$, the solutions of each factor, determined; and, since every sample
gives slightly different factors, which factors of different samples are the same factor?

\paragraph{The angle sets the error bars} On the observation window write $u=\sum_jB_ja_j+\varepsilon=\sumop a+\varepsilon$,
with $B_j$ an orthonormal basis of the sampled kernel of $Q_j$. The noise $\varepsilon\sim\Normal(0,\noise^2I)$ says that each
observation carries independent noise of size $\noise$; the Gaussian prior $a\sim\Normal(0,\omega^2I)$ says that, before
seeing the data, coefficients of size about $\omega$ are expected, and its scale $\omega^2$ is set by empirical Bayes. The posterior of the coefficients is Gaussian,
\begin{equation}\label{eq:posterior-a}
  a\mid u\sim\Normal(\hat a,\Gamma_a),\qquad
  \Gamma_a=\noise^2\Big(\sumop^H\sumop+\frac{\noise^2}{\omega^2}\,I\Big)^{-1};
\end{equation}
the layer returns a distribution of splits, not one split. The covariance has a simple reading: information adds. The
data bring an amount of information $\sumop^H\sumop/\noise^2$ about the coefficients, the prior an amount $1/\omega^2$, and
the posterior variance is one over the total (the derivation is in \SMvariances).

\begin{proposition}[The angle sets the error bars]\label{prop:bayes-angle}
For two groups the eigenvalues of $\sumop^H\sumop$ are $1\pm\cos\theta_k$, with $\theta_k$ the principal angles between the
sampled kernels (and $1$ on the directions of the larger kernel without a partner). In the worst direction orthogonal to
$\ker\sumop$ the posterior variance of the component coefficients is
\begin{equation}\label{eq:bayes-angle}
  \frac{\noise^2}{\gamma_2+\noise^2/\omega^2},\qquad \gamma_2=1-\cF .
\end{equation}
On the shared directions ($\theta_k=0$, the overlap $\deltaP$) the posterior equals the prior, and its mean splits the
shared part equally.
\end{proposition}

Read in the same way, \eqref{eq:bayes-angle} says that in the worst direction the data bring only the information
$\gamma_2/\noise^2$: when the solution spaces of two groups are nearly aligned, the data hardly tell how much of $u$
belongs to each, and the variance grows until the prior caps it. This is the Bayesian form of \eqref{eq:hm}: intervals grow like $\noise/\sin(\thetaF/2)$ as the angle closes and are capped
by the prior when $\thetaF\to0$. The same evidence compares the two-factor model with the one-factor model and gives the
probability $P(\text{split})$, which returns to its prior as $\thetaF\to0$.

\paragraph{The cloud of factors} So far the factors were fixed and the question was about solutions. The second question is
about the factors themselves: every field sample $u^{(q)}$ gave its own list of factors, slightly perturbed from sample to
sample, and the lists are now pooled. Each factor of each sample is mapped to the vector $\beta$ of its monic
coefficients, zero-padded to the largest degrees in play, and enters the cloud with a weight:
\begin{equation}\label{eq:cloud}
  \cloud=\sum_{q=1}^{\Nsamp}\pi_q\sum_jw_j^{(q)}\,\delta_{\beta_j^{(q)}} .
\end{equation}
Here $\pi_q$ is the weight of the field sample $u^{(q)}$, inherited from the hyperparameters it was drawn with
(\cref{sec:modes}), and $w_j^{(q)}$ is the energy share of factor $j$ in that sample, the fraction of the field it explains:
a factor counts in proportion to how plausible its sample is and how much of the field it carries. A sharp equation is a
cloud with a few tight clusters, an uncertain one a diffuse cloud. Since the number of groups is
unknown, the cloud is modelled by a Dirichlet-process Gaussian mixture~\cite{BleJor06} (truncated stick-breaking,
variational inference), which lets the data choose the number of clusters; two clusters that never occur in the same
sample are merged when one Gaussian explains them better. The output is a mixture of Gaussians in coefficient space, one
per group: the centre of each Gaussian is a factor, its spread gives the $90\%$ intervals of the coefficients, its mass
is the share of the cloud it explains, and the cluster responsibilities give the posterior $P(k)$ of the number $k$ of
groups. Because the points of the cloud come from field samples drawn under different hyperparameters, each with its
weight $\pi_q$, the uncertainty about the hyperparameters is part of this spread.
\Cref{fig:grouping} shows the cloud of the running example.

\begin{figure}[tbp]
\centering
\includegraphics[width=\linewidth]{fig_grouping}
\caption{Grouping II on the running example. (a) Each marker is one factor found in one of $32$ posterior samples,
placed by its coefficients of $\xi$ and $\xi^2$; area is proportional to its mass $\pi_qw_j$; crosses mark the true
factors. The cloud has two tight clusters and nothing in between. (b) Largest principal angle between the sampled
kernels of every pair of factors, ordered by cluster: factors of the same group allow the same functions (light
diagonal blocks), factors of different groups do not (dark off-diagonal blocks).}\label{fig:grouping}
\end{figure}

\subsection{Decision and labels}\label{sec:labels}
The interpolated field has already smoothed the high wavenumbers, so the final decisions are made on the \emph{raw
observations}. A candidate set of factors defines a linear model $y=Ea+\varepsilon$ for the $n$ observations, where the
columns of $E$ are the kernel exponentials of the factors evaluated at the observation points. Integrating the
amplitudes out gives
\[
  y\mid Q_1,\dots,Q_N\sim\Normal\big(0,\ \noise^2I+\omega^2EE^H\big):
\]
the data are the sum of the signal $Ea$ and the noise, two independent Gaussians, so their covariances add. The score
\begin{equation}\label{eq:evidence}
  S(Q_1,\dots,Q_N)=\log\Normal\big(y;0,\noise^2I+\omega^2EE^H\big)-\tfrac12n_{\mathrm{par}}\log n
\end{equation}
has two parts. The first is the logarithm of the probability of the observations actually made, under the candidate
factors and with the amplitudes integrated out; the second penalizes the number of coefficients, as in \eqref{eq:bic}.
So $e^{S}$ is, roughly, proportional to the probability that the candidate is the right one among those compared: the
score \emph{compares} candidates and is not a p-value. It decides between dropping a factor, dropping a term whose
interval contains zero, merging two factors or splitting one; a change needs more than one nat. A separate test judges
one candidate on its own: the goodness-of-fit statistic of the observations,
\[
  \chi^2_{\mathrm{obs}}=\sum_{i=1}^n\frac{\abs{y_i-\hat y_i}^2}{\noise^2},
\]
with $\hat y$ the projection of $y$ on the columns of $E$. If the candidate is right it is approximately chi-squared with
$n-\rank E$ degrees of freedom, and its p-value, the probability of a misfit at least this large under the candidate, is
what can reject it. It is distinct from $\chi^2_{\mathrm{modes}}$ of \eqref{eq:bic}, which measures the modes against the
curves, not the observations against the equation. The result then receives one of four labels, which encode
what \cref{sec:certify} allows:

\smallskip
\noindent\begin{tabularx}{\linewidth}{@{}l X@{}}
\toprule
\FC & every factor has more modes than coefficients, at least $d_j+1$ distinct wavenumbers (\cref{prop:counting}), and
the fit passes the goodness-of-fit test: the data are compatible with the factored equation and pin it down within the class.\\
\IDK & the fit is good, but some factor is pinned down only relative to the declared class.\\
\MIS & the goodness-of-fit test rejects the structure ($p<10^{-3}$).\\
\UND & the data cannot decide, typically because of confluent rates or a repeated factor $P_i^2$, whose polynomial
prefactors $t\,e^{\lambda t}$ finite noisy data cannot separate from two nearby rates.\\
\bottomrule
\end{tabularx}
\smallskip

\noindent The output is therefore a distribution, never a point estimate: the factors with $90\%$ intervals, $P(k)$ and
the posterior over structures, the energy shares $w_j$ with intervals, $P(\text{split})$, the separation $\gN$, and the
label. A small Friedrichs angle does not reject the equation; it widens the intervals, so the reader always sees how well
each group is separated, not only whether the equation fits.

